The Length-Width Ratio provides a visually intuitive elongation metric that
complements the perimeter-based and area-based metrics in the K-track compactness
framework.
Our principal findings are:

\begin{enumerate}
  \item The LW $> 3$ court threshold---applied in Illinois and New York
    redistricting litigation---represents a well-established elongation criterion
    that bisect plans satisfy by construction.
  \item All four bisect structure algorithms produce LW $\leq 2.8^\dagger$ on
    NC 2020, with zero districts exceeding the LW $> 3$ threshold across all
    algorithm-state pairs studied.
  \item AABB-based LW can overestimate true LW by up to 41\% for diagonally
    oriented districts; only MBR-based LW (rotating calipers) is
    rotation-invariant and legally defensible.
\end{enumerate}

Practitioners should use MBR-based LW computed via rotating calipers,
explicitly exclude AABB from their methodology, and present LW alongside PP,
Reock, and CH in the composite compactness profile (K.7).
The $\leq 2.8$ maximum LW for all bisect NC plans provides a strong
certified result for litigation contexts where elongation is at issue.
